%\documentclass[spanish,a4paper,12pt,oneside]{extreport}

\documentclass[12pt, a4paper]{report}

\pagestyle{empty}

\usepackage{pxfonts}
\usepackage[T1]{fontenc}

\usepackage{booktabs}

\usepackage{graphicx}

\usepackage[linesnumbered, ruled, vlined, spanish]{algorithm2e}

\usepackage{hyperref}

%\usepackage{algorithm}
\usepackage{algpseudocode}

\usepackage{tikz}

\usetikzlibrary{shapes.geometric, arrows.meta, positioning, calc, backgrounds}

% Definición de colores basados en la imagen
\definecolor{brown_main}{RGB}{139, 69, 19}
\definecolor{blue_main}{RGB}{30, 50, 100}
\definecolor{bg_green}{RGB}{245, 250, 245}
\definecolor{bg_blue}{RGB}{245, 247, 255}
\definecolor{node_yellow}{RGB}{255, 252, 235}
\definecolor{border_yellow}{RGB}{218, 165, 32}

\usepackage{multirow}
\usepackage{rotating}
\usepackage{amssymb}
\usepackage{amsbsy}
\usepackage{amsthm}
\usepackage{arydshln}
\usepackage{color}
\usepackage{amsmath}

\usepackage[backend=biber,style=authoryear,citestyle=apa]{biblatex}
\addbibresource{Biblio.bib}

\usepackage{lineno}
\usepackage[utf8]{inputenc}
\usepackage[spanish]{babel}

% Cambiar "Cuadro" por "Tabla"
\addto\captionsspanish{\renewcommand{\tablename}{Tabla}}
\addto\captionsspanish{\renewcommand{\listtablename}{Índice de tablas}}

\usetikzlibrary{shapes}

\begin{document}

\renewcommand\listtablename{Índice de Tablas}
\renewcommand\listfigurename{Índice de Figuras}

\pagestyle{empty}
\thispagestyle{empty}

\newcommand{\HRule}{\rule{\linewidth}{1mm}}
\setlength{\parindent}{0mm}
\setlength{\parskip}{0mm}

\vspace*{\stretch{0.5}}

\begin{center}
\includegraphics[scale=0.8]{images/escuela-ingenieria-tecnologia-original}\\[10mm]
{\Huge Informe de la Práctica 5 de DAA \\[3mm] Curso 2025-2026}
\end{center}

\HRule
\begin{center}
        {\Huge \itshape Multi-Source \\ Capacitated Facility Location Problem \\ with Customer Incompatibilities \\ \vspace{10pt} (MS-CFLP-CI)} \\[5mm]
        {\Large Kyliam Chinea Salcedo} \\[5mm]
\end{center}
\HRule
\vspace*{\stretch{2}}
\begin{center}
  \Large La Laguna, \today
\end{center}

\setlength{\parindent}{5mm}

\newpage

\tableofcontents
\newpage
\listoffigures
\newpage
\listoftables

% Numeración a partir del capítulo I
\renewcommand{\thepage}{\arabic{page}}
\setcounter{page}{1}
\pagestyle{plain}

\chapter{Introducción}

Esta práctica aborda el \emph{Multi-Source Capacitated Facility Location Problem with Customer Incompatibilities} (MS-CFLP-CI) usando una combinación de:

\begin{itemize}
  \item Constructivas (Voraz y GRASP).
  \item Búsqueda local basada en movimientos (\emph{moves}) con cálculo de $\Delta$-coste.
  \item Metaheurísticas (RVND y GVNS) para escapar de óptimos locales.
\end{itemize}

El objetivo principal es minimizar el coste total combinando costes fijos de apertura de instalaciones y costes variables de transporte, cumpliendo además:

\begin{itemize}
  \item Restricciones de capacidad.
  \item Satisfacción completa de demanda (multi-source).
  \item Incompatibilidades entre clientes servidos desde una misma instalación.
\end{itemize}


\chapter{MS-CFLP-CI: Modelo y restricciones}

\section{Conjuntos y parámetros}

Sea $I$ el conjunto de clientes y $J$ el conjunto de instalaciones.

\begin{itemize}
  \item $d_i$ demanda del cliente $i\in I$.
  \item $cap_j$ capacidad de la instalación $j\in J$.
  \item $f_j$ coste fijo de abrir $j$.
  \item $c_{ij}$ coste unitario de transporte desde $j$ hasta $i$.
  \item $Inc\subseteq I\times I$ conjunto de pares incompatibles.
\end{itemize}

\section{Variables de decisión}

\begin{itemize}
  \item $x_{ij}\ge 0$: cantidad servida al cliente $i$ desde la instalación $j$.
  \item $y_j\in\{0,1\}$: vale 1 si la instalación $j$ está abierta.
\end{itemize}

\section{Función objetivo}

\[
\min \sum_{j\in J} f_j y_j + \sum_{i\in I} \sum_{j\in J} c_{ij} x_{ij}
\]

\section{Restricciones}

\subsection{Demanda (multi-source)}

Cada cliente debe recibir exactamente su demanda total:

\[
\sum_{j\in J} x_{ij} = d_i \quad \forall i\in I.
\]

\subsection{Capacidad de instalaciones}

\[
\sum_{i\in I} x_{ij} \le cap_j \quad \forall j\in J.
\]

\subsection{Incompatibilidades}

Si $ (i,k)\in Inc$, entonces no pueden coexistir en la misma instalación:

\[
\forall j\in J: (x_{ij}>0)\land(x_{kj}>0) \;\Rightarrow\; (i,k)\notin Inc.
\]


\section{Representación de la solución en la implementación}

En el código, la solución se representa con una estructura incremental (clase \texttt{model.Solucion}) basada en:

\begin{itemize}
  \item Una matriz $x$ almacenada como \texttt{double[][] suministros} donde $x_{ij}$ es la cantidad que la instalación $j$ suministra al cliente $i$.
  \item Un vector \texttt{capacidadRestante[j]} que mantiene $cap_j - \sum_i x_{ij}$.
  \item Un vector booleano \texttt{instalacionesAbiertas[j]} para el estado $y_j$.
  \item Un coste total \texttt{costeTotal} actualizado en cada modificación.
  \item Un \texttt{BitSet} por instalación con los clientes presentes (para compatibilidad rápida).
\end{itemize}

\subsection{Por qué se escogió esta representación (y por qué funciona)}

\begin{enumerate}
  \item \textbf{Permite multi-source de forma natural:} como $x_{ij}$ es real no negativa, un cliente puede dividir su demanda entre varias instalaciones.

  \item \textbf{Soporta movimientos locales eficientemente:} los moves cambian pocos elementos de $x$. Si el coste se recalcula desde cero en cada intento, el tiempo se dispara.

  \item \textbf{Coste incremental (clave):} al modificar $x_{ij}$ en una cantidad $q$:
  \[
    \Delta_{var} = q\,c_{ij}
  \]
  y por tanto el coste variable se actualiza con una suma/resta local.

  \item \textbf{Gestión automática del coste fijo:} abrir/cerrar una instalación es una condición local: si una instalación pasa de $0$ carga a $>0$ carga, se abre y suma $f_j$; si se queda sin carga, se cierra y resta $f_j$.

  \item \textbf{Compatibilidad rápida:} al almacenar el conjunto de clientes por instalación en bitsets, comprobar si un cliente es compatible con lo ya presente se reduce a una intersección de bitsets.
\end{enumerate}

\subsection{Cómo impacta esta representación en los algoritmos}

\begin{itemize}
  \item En \textbf{Voraz} y \textbf{GRASP}, la capacidad restante y la compatibilidad se consultan constantemente; tenerlas precomputadas acelera cada asignación.
  \item En \textbf{Búsqueda local}, la evaluación de un movimiento depende de deltas. El coste incremental evita tener que recalcular la función objetivo completa.
  \item En \textbf{SwapInstalaciones}, simular movimientos grandes requiere aplicar muchas operaciones y después deshacerlas: el rollback es barato porque las operaciones elementales ya actualizan todas las estructuras auxiliares (capacidad, bitsets, coste).
\end{itemize}


\chapter{Algoritmos y búsqueda local (paso a paso)}

\section{Búsqueda local basada en \emph{moves}}

La implementación se basa en dos interfaces:

\begin{itemize}
  \item \texttt{BusquedaLocal}: una vecindad que sabe encontrar el mejor movimiento factible.
  \item \texttt{Move}: un objeto que sabe calcular $\Delta$-coste y aplicarse.
\end{itemize}

\subsection{Esquema general}

Dada una solución actual $S$:

\begin{enumerate}
  \item Una vecindad explora candidatos y devuelve el mejor move $m$ (si existe).
  \item Se calcula $\Delta = f(S\oplus m) - f(S)$.
  \item Si $\Delta < -\varepsilon$ se aplica el movimiento.
\end{enumerate}

\subsection{Por qué el $\Delta$-coste es correcto}

La función objetivo es suma de costes fijos y variables:

\[
 f(S) = \sum_{j\in J} f_j y_j + \sum_{i\in I} \sum_{j\in J} c_{ij} x_{ij}.
\]

Un move sólo modifica unas pocas asignaciones $x_{ij}$ y, a lo sumo, el estado de algunas $y_j$. Por eso $\Delta$ se puede calcular usando sólo los términos afectados.

\subsection{Ejemplo: ShiftMove}

Mover una cantidad $q$ del cliente $i$ desde $j_1$ a $j_2$ implica:

\[
\Delta_{var} = q\,(c_{i j_2} - c_{i j_1}).
\]

Además:

\begin{itemize}
  \item Si $j_2$ estaba cerrada y el move la abre: $\Delta_{fijo} += f_{j_2}$.
  \item Si tras quitar $q$ la instalación $j_1$ queda vacía y se cierra: $\Delta_{fijo} -= f_{j_1}$.
\end{itemize}

Por tanto:

\[
\Delta = \Delta_{var} + \Delta_{fijo}.
\]

\subsection{Moves compuestos y rollback: SwapInstalaciones}

\textbf{Idea:} cerrar una instalación abierta $j_{open}$ y reubicar sus cargas (posiblemente abriendo una instalación cerrada $j_{closed}$).

\textbf{Por qué hace falta rollback:} para evaluar un candidato, hay que simular muchas operaciones de quitar/añadir suministro. Si se deja el estado modificado, la evaluación de otros candidatos se corrompe.

\textbf{Estrategia implementada:}

\begin{enumerate}
  \item Se aplica una secuencia de operaciones atómicas (quitar/añadir) sobre la solución.
  \item Se registra cada operación.
  \item Si el candidato no es factible, o tras calcular el delta, se deshace todo en orden inverso.
\end{enumerate}

Si una operación fue “add”, su inversa es “remove”, y viceversa. Este rollback es correcto porque las operaciones elementales actualizan de manera consistente coste, capacidad y bitsets.

\section{RVND: Randomized Variable Neighborhood Descent}

RVND alterna vecindades en orden aleatorio para evitar sesgos.

\begin{algorithm}[H]
\DontPrintSemicolon
\caption{RVND (esquema implementado)}
\KwIn{Solución inicial $S_0$, conjunto de vecindades $\mathcal{N}$}
\KwOut{Mínimo local $S$}

$S \leftarrow S_0$; $D \leftarrow \mathcal{N}$\;
\While{$D$ no vacío}{
  elegir $N$ aleatoriamente de $D$\;
  $S' \leftarrow$ aplicar como máximo un move de $N$ sobre una copia de $S$\;
  \If{$f(S') < f(S) - \varepsilon$}{
    $S \leftarrow S'$; $D \leftarrow \mathcal{N}$\;
  }{
    eliminar $N$ de $D$\;
  }
}
\Return{$S$}
\end{algorithm}

\section{GVNS-RL: GVNS con ruleta ponderada (reinforcement learning)}

GVNS alterna:

\begin{itemize}
  \item \textbf{Shaking}: perturbar la mejor solución actual eliminando y reconstruyendo parte de la asignación.
  \item \textbf{Descenso}: aplicar un RVND interno, donde la selección de vecindad se hace por ruleta según pesos.
\end{itemize}

\subsection{Bucle principal}

\begin{enumerate}
  \item Inicializar $S_{best} \leftarrow S_0$, $k \leftarrow 1$.
  \item Mientras no se alcance el límite de estancamiento:
  \begin{enumerate}
    \item $S_{shake} \leftarrow$ perturbar($S_{best}$, $k$).
    \item $S_{loc} \leftarrow$ RVND-RL($S_{shake}$).
    \item Si $f(S_{loc}) < f(S_{best}) - \varepsilon$, aceptar y reiniciar $k\leftarrow 1$.
    \item Si no mejora, aumentar $k$ hasta $k_{max}$ (cíclico).
  \end{enumerate}
\end{enumerate}

\subsection{Ruleta ponderada (RL) en el RVND interno}

Cada vecindad $N_i$ tiene un peso $w_i$ (inicialmente 1). La ruleta elige $N_i$ con probabilidad:

\[
P(N_i) = \frac{w_i}{\sum_{h \in D} w_h}.
\]

Si $N_i$ mejora, se refuerza ($w_i \leftarrow w_i + 1$). Si no mejora, se penaliza ($w_i \leftarrow \max(0.1, 0.9\,w_i)$).


\chapter{Experimentos y resultados computacionales}

\section{Tablas exportadas por el programa}

Para cada instancia se exportan ficheros \texttt{Tabla*.txt} en \texttt{tablas/<instancia>/}. Para resultados comparativos usamos:

\begin{itemize}
  \item \texttt{Tabla10.txt}: resultados de GRASP (incluye fila \emph{Promedio}).
  \item \texttt{Tabla11.txt}: resultados de GVNS (incluye fila \emph{Promedio}).
  \item \texttt{Tabla8.txt}: desglose de costes de la mejor solución final exportada.
\end{itemize}

\section{Resumen wlp01--wlp05}

A partir de las tablas exportadas, calculamos:

\begin{itemize}
  \item Promedio de coste total y tiempo para GRASP.
  \item Promedio de coste total y tiempo para GVNS.
  \item Mejor coste final (Tabla 8).
  \item Mejora porcentual promedio de GVNS respecto a GRASP.
\end{itemize}

\begin{table}[H]
\centering
\caption{Resumen comparativo (promedios) para wlp01--wlp05}
\label{tab:resumen}
\begin{tabular}{lrrrrrr}
\toprule
Instancia & GRASP $\overline{C}$ & GRASP $\overline{t}$ (s) & GVNS $\overline{C}$ & GVNS $\overline{t}$ (s) & Mejor $C$ & Mejora \\ 
\midrule
wlp01 & 39147.33 & 0.06 & 30064.50 & 0.29 & 30191 & 23.20\% \\ 
wlp02 & 73774.33 & 0.11 & 57242.00 & 2.98 & 55835 & 22.41\% \\ 
wlp03 & 89114.33 & 0.25 & 68323.33 & 8.71 & 70308 & 23.33\% \\ 
wlp04 & 118318.33 & 0.73 & 89521.17 & 31.78 & 93540 & 24.34\% \\ 
wlp05 & 146666.00 & 2.05 & 111828.17 & 105.19 & 111406 & 23.75\% \\ 
\bottomrule
\end{tabular}
\end{table}

\begin{table}[H]
\centering
\caption{Desglose del mejor coste final (Tabla 8) para wlp01--wlp05}
\label{tab:costes}
\begin{tabular}{lrrr}
\toprule
Instancia & Coste fijo & Coste variable & Coste total \\ 
\midrule
wlp01 & 16050 & 14141 & 30191 \\ 
wlp02 & 32430 & 23405 & 55835 \\ 
wlp03 & 42850 & 27458 & 70308 \\ 
wlp04 & 57500 & 36040 & 93540 \\ 
wlp05 & 74860 & 36546 & 111406 \\ 
\bottomrule
\end{tabular}
\end{table}

\section{Análisis}

Se observa una reducción sistemática del coste total al pasar de GRASP a GVNS, a cambio de un incremento notable del tiempo de CPU.
Este comportamiento es consistente con la lógica de GVNS:

\begin{itemize}
  \item El \emph{shaking} explora regiones nuevas del espacio de soluciones.
  \item El descenso (RVND-RL) intensifica alrededor de la región alcanzada.
  \item El uso de pesos en la ruleta favorece vecindades que históricamente han mejorado.
\end{itemize}


\chapter{Conclusiones y trabajo futuro}

\begin{itemize}
  \item La representación incremental de la solución permite calcular $\Delta$-costes y hacer rollback eficientemente.
  \item La separación \emph{vecindad} vs \emph{move} hace el código modular y fácil de extender.
  \item GVNS-RL consigue mejoras fuertes sobre GRASP en las instancias estudiadas.
  \item Como trabajo futuro, se puede reforzar la fase de \emph{shaking} para evitar reconstrucciones parciales y explorar estrategias de aceptación más sofisticadas.
\end{itemize}

\printbibliography[heading=bibintoc, title={Bibliografía}]
\end{document}
